% --- Archon physics formalization source begin ---
% archon:physics
% archon:covers PhyXMiniProblems/problem_phyx_mini_0364.lean
% archon:source-report reports/phyx_mini/problem_phyx_mini_0364.source.json

\chapter{Physics problem phyx\_mini\_0364}
\label{ch:PhyXMiniProblems_problem_phyx_mini_0364}

\paragraph{Problem source.}
\#\# Physical scenario

The 3.0-m-long pipe in the figure is closed at the top end. It is slowly pushed straight down into the water until the top end of the pipe is level with the water’s surface.

\#\# Question

What is the length of the trapped air?

\#\# Answer choices

- A: A: 2.4m

- B: B: 2.2m

- C: C: 2.0m

- D: D: 2.1m

\#\# Figure caption (auxiliary; use the image as primary evidence)

The image shows two diagrams side by side, labeled "Before" and "After." Both diagrams feature a vertical rectangular object. 

1. **"Before" Diagram:**
   - The rectangular object is fully above a blue surface.
   - It has an annotated height of "3.0 m."
   - The background is white, indicating it's not submerged.

2. **"After" Diagram:**
   - The rectangular object is partially submerged in a blue area representing a liquid.
   - The portion above the liquid is labeled with "L" indicating a new height above the liquid.
   - The background inside the liquid is blue.

The diagrams illustrate the change in the position of the object relative to a liquid surface.

\#\# Dataset metadata

Ideal Gases and Kinetic Theory; Physical Model Grounding Reasoning, Multi-Formula Reasoning

\paragraph{Recorded answer/context.}
A: A: 2.4m

\paragraph{Figure/image path.}
/root/proposal\_for\_physic/science-mango/phyx\_mini\_run/phyx\_data/test\_image/364.png

\paragraph{Formalization target.}
create a compiling Lean file with sorry bodies at `PhyXMiniProblems/problem\_phyx\_mini\_0364.lean`. The Lean declarations must preserve the physical quantities, dimensions or dimensional roles, figure labels, governing-law hypotheses, and final relation expressed by this problem.
Use Mathlib/Physlib names found through LeanExplore where available. If a physics API is missing, introduce faithful local abstractions rather than scalar placeholder aliases.

\begin{theorem}[Physics formalization target]
\label{thm:physics:phyx_mini_0364:target}
\lean{PhyXMiniProblems.ProblemPhyXMini0364.problem_phyx_mini_0364}
\uses{def:physics:phyx-mini-0364:phyxminiproblems-problemphyxmini0364-lengthinmeters, def:physics:phyx-mini-0364:phyxminiproblems-problemphyxmini0364-closedtoppipesetup, def:physics:phyx-mini-0364:phyxminiproblems-problemphyxmini0364-matchesprimaryfigure, def:physics:phyx-mini-0364:phyxminiproblems-problemphyxmini0364-matchesclosedtoppipescenario, def:physics:phyx-mini-0364:phyxminiproblems-problemphyxmini0364-usesstandardatmospherewaterandgravity, def:physics:phyx-mini-0364:phyxminiproblems-problemphyxmini0364-hasphysicalpipeparameters, def:physics:phyx-mini-0364:phyxminiproblems-problemphyxmini0364-satisfiesclosedtoppipelaws, def:physics:phyx-mini-0364:phyxminiproblems-problemphyxmini0364-displayedlengthinmeters, def:physics:phyx-mini-0364:phyxminiproblems-problemphyxmini0364-roundstonearesttenthmeter, def:physics:phyx-mini-0364:phyxminiproblems-problemphyxmini0364-isuniquematchingdisplayedlength}
The assigned autoformalize agent should translate this physics problem into Lean declarations in the covered file, with theorem and lemma proof bodies written as `by sorry`.
\end{theorem}
\begin{proof}
This is an autoformalization task, not a proof task. Produce statements that can later be proved by the physics prover without weakening the physical model.
\end{proof}

% --- Archon declaration topology begin ---
\section{Declaration topology}
\begin{definition}[volume Dimension]
  \label{def:physics:phyx-mini-0364:phyxminiproblems-problemphyxmini0364-volumedimension}
  \lean{PhyXMiniProblems.ProblemPhyXMini0364.volumeDimension}
  The physical dimension of volume, L\textasciicircum{}3
\end{definition}
\begin{proof}
  This is the declared model definition of volume Dimension.
\end{proof}
\begin{definition}[mass Density Dimension]
  \label{def:physics:phyx-mini-0364:phyxminiproblems-problemphyxmini0364-massdensitydimension}
  \lean{PhyXMiniProblems.ProblemPhyXMini0364.massDensityDimension}
  The physical dimension of mass density, M L\textasciicircum{}-3
\end{definition}
\begin{proof}
  This is the declared model definition of mass Density Dimension.
\end{proof}
\begin{definition}[acceleration Dimension]
  \label{def:physics:phyx-mini-0364:phyxminiproblems-problemphyxmini0364-accelerationdimension}
  \lean{PhyXMiniProblems.ProblemPhyXMini0364.accelerationDimension}
  The physical dimension of acceleration, L T\textasciicircum{}-2
\end{definition}
\begin{proof}
  This is the declared model definition of acceleration Dimension.
\end{proof}
\begin{definition}[Length Quantity]
  \label{def:physics:phyx-mini-0364:phyxminiproblems-problemphyxmini0364-lengthquantity}
  \lean{PhyXMiniProblems.ProblemPhyXMini0364.LengthQuantity}
  A nonnegative physical length, independent of the readout unit
\end{definition}
\begin{proof}
  This is the declared model definition of Length Quantity.
\end{proof}
\begin{definition}[Area Quantity]
  \label{def:physics:phyx-mini-0364:phyxminiproblems-problemphyxmini0364-areaquantity}
  \lean{PhyXMiniProblems.ProblemPhyXMini0364.AreaQuantity}
  A nonnegative physical cross-sectional area
\end{definition}
\begin{proof}
  This is the declared model definition of Area Quantity.
\end{proof}
\begin{definition}[Volume Quantity]
  \label{def:physics:phyx-mini-0364:phyxminiproblems-problemphyxmini0364-volumequantity}
  \lean{PhyXMiniProblems.ProblemPhyXMini0364.VolumeQuantity}
  \uses{def:physics:phyx-mini-0364:phyxminiproblems-problemphyxmini0364-volumedimension}
  A nonnegative physical volume
\end{definition}
\begin{proof}
  This is the declared model definition of Volume Quantity.
\end{proof}
\begin{definition}[Mass Density Quantity]
  \label{def:physics:phyx-mini-0364:phyxminiproblems-problemphyxmini0364-massdensityquantity}
  \lean{PhyXMiniProblems.ProblemPhyXMini0364.MassDensityQuantity}
  \uses{def:physics:phyx-mini-0364:phyxminiproblems-problemphyxmini0364-massdensitydimension}
  A nonnegative physical mass density
\end{definition}
\begin{proof}
  This is the declared model definition of Mass Density Quantity.
\end{proof}
\begin{definition}[Acceleration Quantity]
  \label{def:physics:phyx-mini-0364:phyxminiproblems-problemphyxmini0364-accelerationquantity}
  \lean{PhyXMiniProblems.ProblemPhyXMini0364.AccelerationQuantity}
  \uses{def:physics:phyx-mini-0364:phyxminiproblems-problemphyxmini0364-accelerationdimension}
  A nonnegative physical acceleration magnitude
\end{definition}
\begin{proof}
  This is the declared model definition of Acceleration Quantity.
\end{proof}
\begin{definition}[length Readout]
  \label{def:physics:phyx-mini-0364:phyxminiproblems-problemphyxmini0364-lengthreadout}
  \lean{PhyXMiniProblems.ProblemPhyXMini0364.lengthReadout}
  \uses{def:physics:phyx-mini-0364:phyxminiproblems-problemphyxmini0364-lengthquantity}
  Read a physical length in a selected length unit
\end{definition}
\begin{proof}
  This is the declared model definition of length Readout.
\end{proof}
\begin{definition}[area Readout]
  \label{def:physics:phyx-mini-0364:phyxminiproblems-problemphyxmini0364-areareadout}
  \lean{PhyXMiniProblems.ProblemPhyXMini0364.areaReadout}
  \uses{def:physics:phyx-mini-0364:phyxminiproblems-problemphyxmini0364-areaquantity}
  Read a physical area in the square of a selected length unit
\end{definition}
\begin{proof}
  This is the declared model definition of area Readout.
\end{proof}
\begin{definition}[volume Readout]
  \label{def:physics:phyx-mini-0364:phyxminiproblems-problemphyxmini0364-volumereadout}
  \lean{PhyXMiniProblems.ProblemPhyXMini0364.volumeReadout}
  \uses{def:physics:phyx-mini-0364:phyxminiproblems-problemphyxmini0364-volumequantity}
  Read a physical volume in the cube of a selected length unit
\end{definition}
\begin{proof}
  This is the declared model definition of volume Readout.
\end{proof}
\begin{definition}[pressure Readout]
  \label{def:physics:phyx-mini-0364:phyxminiproblems-problemphyxmini0364-pressurereadout}
  \lean{PhyXMiniProblems.ProblemPhyXMini0364.pressureReadout}
  Read pressure in the coherent unit induced by selected base units
\end{definition}
\begin{proof}
  This is the declared model definition of pressure Readout.
\end{proof}
\begin{definition}[density Readout]
  \label{def:physics:phyx-mini-0364:phyxminiproblems-problemphyxmini0364-densityreadout}
  \lean{PhyXMiniProblems.ProblemPhyXMini0364.densityReadout}
  \uses{def:physics:phyx-mini-0364:phyxminiproblems-problemphyxmini0364-massdensityquantity}
  Read mass density in a selected mass unit per selected length unit cubed
\end{definition}
\begin{proof}
  This is the declared model definition of density Readout.
\end{proof}
\begin{definition}[acceleration Readout]
  \label{def:physics:phyx-mini-0364:phyxminiproblems-problemphyxmini0364-accelerationreadout}
  \lean{PhyXMiniProblems.ProblemPhyXMini0364.accelerationReadout}
  \uses{def:physics:phyx-mini-0364:phyxminiproblems-problemphyxmini0364-accelerationquantity}
  Read acceleration in selected length and time units
\end{definition}
\begin{proof}
  This is the declared model definition of acceleration Readout.
\end{proof}
\begin{definition}[length In Meters]
  \label{def:physics:phyx-mini-0364:phyxminiproblems-problemphyxmini0364-lengthinmeters}
  \lean{PhyXMiniProblems.ProblemPhyXMini0364.lengthInMeters}
  \uses{def:physics:phyx-mini-0364:phyxminiproblems-problemphyxmini0364-lengthquantity, def:physics:phyx-mini-0364:phyxminiproblems-problemphyxmini0364-lengthreadout}
  Meter readout of a physical length
\end{definition}
\begin{proof}
  This is the declared model definition of length In Meters.
\end{proof}
\begin{definition}[area In Square Meters]
  \label{def:physics:phyx-mini-0364:phyxminiproblems-problemphyxmini0364-areainsquaremeters}
  \lean{PhyXMiniProblems.ProblemPhyXMini0364.areaInSquareMeters}
  \uses{def:physics:phyx-mini-0364:phyxminiproblems-problemphyxmini0364-areaquantity, def:physics:phyx-mini-0364:phyxminiproblems-problemphyxmini0364-areareadout}
  Square-meter readout of a physical area
\end{definition}
\begin{proof}
  This is the declared model definition of area In Square Meters.
\end{proof}
\begin{definition}[volume In Cubic Meters]
  \label{def:physics:phyx-mini-0364:phyxminiproblems-problemphyxmini0364-volumeincubicmeters}
  \lean{PhyXMiniProblems.ProblemPhyXMini0364.volumeInCubicMeters}
  \uses{def:physics:phyx-mini-0364:phyxminiproblems-problemphyxmini0364-volumequantity, def:physics:phyx-mini-0364:phyxminiproblems-problemphyxmini0364-volumereadout}
  Cubic-meter readout of a physical volume
\end{definition}
\begin{proof}
  This is the declared model definition of volume In Cubic Meters.
\end{proof}
\begin{definition}[pressure In Pascals]
  \label{def:physics:phyx-mini-0364:phyxminiproblems-problemphyxmini0364-pressureinpascals}
  \lean{PhyXMiniProblems.ProblemPhyXMini0364.pressureInPascals}
  \uses{def:physics:phyx-mini-0364:phyxminiproblems-problemphyxmini0364-pressurereadout}
  Pascal readout of an absolute pressure
\end{definition}
\begin{proof}
  This is the declared model definition of pressure In Pascals.
\end{proof}
\begin{definition}[density In Kilograms Per Cubic Meter]
  \label{def:physics:phyx-mini-0364:phyxminiproblems-problemphyxmini0364-densityinkilogramspercubicmeter}
  \lean{PhyXMiniProblems.ProblemPhyXMini0364.densityInKilogramsPerCubicMeter}
  \uses{def:physics:phyx-mini-0364:phyxminiproblems-problemphyxmini0364-massdensityquantity, def:physics:phyx-mini-0364:phyxminiproblems-problemphyxmini0364-densityreadout}
  Kilogram-per-cubic-meter readout of a mass density
\end{definition}
\begin{proof}
  This is the declared model definition of density In Kilograms Per Cubic Meter.
\end{proof}
\begin{definition}[acceleration In Meters Per Second Squared]
  \label{def:physics:phyx-mini-0364:phyxminiproblems-problemphyxmini0364-accelerationinmeterspersecondsquared}
  \lean{PhyXMiniProblems.ProblemPhyXMini0364.accelerationInMetersPerSecondSquared}
  \uses{def:physics:phyx-mini-0364:phyxminiproblems-problemphyxmini0364-accelerationquantity, def:physics:phyx-mini-0364:phyxminiproblems-problemphyxmini0364-accelerationreadout}
  Meter-per-second-squared readout of an acceleration
\end{definition}
\begin{proof}
  This is the declared model definition of acceleration In Meters Per Second Squared.
\end{proof}
\begin{definition}[Pipe State]
  \label{def:physics:phyx-mini-0364:phyxminiproblems-problemphyxmini0364-pipestate}
  \lean{PhyXMiniProblems.ProblemPhyXMini0364.PipeState}
  The two physical states pictured in the supplied before/after figure
\end{definition}
\begin{proof}
  This is the declared model definition of Pipe State.
\end{proof}
\begin{definition}[Pipe Liquid]
  \label{def:physics:phyx-mini-0364:phyxminiproblems-problemphyxmini0364-pipeliquid}
  \lean{PhyXMiniProblems.ProblemPhyXMini0364.PipeLiquid}
  The liquid surrounding and entering the pipe
\end{definition}
\begin{proof}
  This is the declared model definition of Pipe Liquid.
\end{proof}
\begin{definition}[Pipe Gas]
  \label{def:physics:phyx-mini-0364:phyxminiproblems-problemphyxmini0364-pipegas}
  \lean{PhyXMiniProblems.ProblemPhyXMini0364.PipeGas}
  The gas occupying the pipe
\end{definition}
\begin{proof}
  This is the declared model definition of Pipe Gas.
\end{proof}
\begin{definition}[Immersion Protocol]
  \label{def:physics:phyx-mini-0364:phyxminiproblems-problemphyxmini0364-immersionprotocol}
  \lean{PhyXMiniProblems.ProblemPhyXMini0364.ImmersionProtocol}
  How the pipe is moved from the before state to the after state
\end{definition}
\begin{proof}
  This is the declared model definition of Immersion Protocol.
\end{proof}
\begin{definition}[Compression Protocol]
  \label{def:physics:phyx-mini-0364:phyxminiproblems-problemphyxmini0364-compressionprotocol}
  \lean{PhyXMiniProblems.ProblemPhyXMini0364.CompressionProtocol}
  Thermal model for the slowly compressed trapped air
\end{definition}
\begin{proof}
  This is the declared model definition of Compression Protocol.
\end{proof}
\begin{definition}[Figure Text Label]
  \label{def:physics:phyx-mini-0364:phyxminiproblems-problemphyxmini0364-figuretextlabel}
  \lean{PhyXMiniProblems.ProblemPhyXMini0364.FigureTextLabel}
  Literal labels visible in the supplied bitmap
\end{definition}
\begin{proof}
  This is the declared model definition of Figure Text Label.
\end{proof}
\begin{definition}[Figure Length Mark]
  \label{def:physics:phyx-mini-0364:phyxminiproblems-problemphyxmini0364-figurelengthmark}
  \lean{PhyXMiniProblems.ProblemPhyXMini0364.FigureLengthMark}
  The two dimensionful length marks drawn in the figure
\end{definition}
\begin{proof}
  This is the declared model definition of Figure Length Mark.
\end{proof}
\begin{definition}[Figure Length Role]
  \label{def:physics:phyx-mini-0364:phyxminiproblems-problemphyxmini0364-figurelengthrole}
  \lean{PhyXMiniProblems.ProblemPhyXMini0364.FigureLengthRole}
  Physical roles of the length marks in the supplied figure
\end{definition}
\begin{proof}
  This is the declared model definition of Figure Length Role.
\end{proof}
\begin{definition}[Closed Top Pipe Figure]
  \label{def:physics:phyx-mini-0364:phyxminiproblems-problemphyxmini0364-closedtoppipefigure}
  \lean{PhyXMiniProblems.ProblemPhyXMini0364.ClosedTopPipeFigure}
  \uses{def:physics:phyx-mini-0364:phyxminiproblems-problemphyxmini0364-lengthquantity, def:physics:phyx-mini-0364:phyxminiproblems-problemphyxmini0364-figuretextlabel, def:physics:phyx-mini-0364:phyxminiproblems-problemphyxmini0364-figurelengthmark, def:physics:phyx-mini-0364:phyxminiproblems-problemphyxmini0364-figurelengthrole}
  Dimensionful and qualitative data represented by the primary image. The L arrow has no numerical value stored here; it is an independently modeled physical length that the theorem must determine
\end{definition}
\begin{proof}
  This is the declared model definition of Closed Top Pipe Figure.
\end{proof}
\begin{definition}[Closed Top Pipe Setup]
  \label{def:physics:phyx-mini-0364:phyxminiproblems-problemphyxmini0364-closedtoppipesetup}
  \lean{PhyXMiniProblems.ProblemPhyXMini0364.ClosedTopPipeSetup}
  \uses{def:physics:phyx-mini-0364:phyxminiproblems-problemphyxmini0364-lengthquantity, def:physics:phyx-mini-0364:phyxminiproblems-problemphyxmini0364-areaquantity, def:physics:phyx-mini-0364:phyxminiproblems-problemphyxmini0364-volumequantity, def:physics:phyx-mini-0364:phyxminiproblems-problemphyxmini0364-massdensityquantity, def:physics:phyx-mini-0364:phyxminiproblems-problemphyxmini0364-accelerationquantity, def:physics:phyx-mini-0364:phyxminiproblems-problemphyxmini0364-pipestate, def:physics:phyx-mini-0364:phyxminiproblems-problemphyxmini0364-pipeliquid, def:physics:phyx-mini-0364:phyxminiproblems-problemphyxmini0364-pipegas, def:physics:phyx-mini-0364:phyxminiproblems-problemphyxmini0364-immersionprotocol, def:physics:phyx-mini-0364:phyxminiproblems-problemphyxmini0364-compressionprotocol, def:physics:phyx-mini-0364:phyxminiproblems-problemphyxmini0364-closedtoppipefigure}
  Independent physical quantities in the immersion experiment. In particular, the final trapped-air length and pressure are fields rather than definitions made from the recorded answer or from the desired closed form
\end{definition}
\begin{proof}
  This is the declared model definition of Closed Top Pipe Setup.
\end{proof}
\begin{definition}[Matches Primary Figure]
  \label{def:physics:phyx-mini-0364:phyxminiproblems-problemphyxmini0364-matchesprimaryfigure}
  \lean{PhyXMiniProblems.ProblemPhyXMini0364.MatchesPrimaryFigure}
  \uses{def:physics:phyx-mini-0364:phyxminiproblems-problemphyxmini0364-lengthinmeters, def:physics:phyx-mini-0364:phyxminiproblems-problemphyxmini0364-closedtoppipesetup}
  Evidence read from the primary bitmap. It records both named panels, the 3.0 m mark, the role of L, and the visible water/air geometry. It does not assign a numerical value to L
\end{definition}
\begin{proof}
  This is the declared model definition of Matches Primary Figure.
\end{proof}
\begin{definition}[Matches Closed Top Pipe Scenario]
  \label{def:physics:phyx-mini-0364:phyxminiproblems-problemphyxmini0364-matchesclosedtoppipescenario}
  \lean{PhyXMiniProblems.ProblemPhyXMini0364.MatchesClosedTopPipeScenario}
  \uses{def:physics:phyx-mini-0364:phyxminiproblems-problemphyxmini0364-closedtoppipesetup}
  Qualitative prose and the vertical geometry of the experiment. The initial air fills the pipe. In the final state the closed top is at the outside water surface, so the depth of the internal interface equals the final air-column length. These relations do not specify that length numerically
\end{definition}
\begin{proof}
  This is the declared model definition of Matches Closed Top Pipe Scenario.
\end{proof}
\begin{definition}[Uses Standard Atmosphere Water And Gravity]
  \label{def:physics:phyx-mini-0364:phyxminiproblems-problemphyxmini0364-usesstandardatmospherewaterandgravity}
  \lean{PhyXMiniProblems.ProblemPhyXMini0364.UsesStandardAtmosphereWaterAndGravity}
  \uses{def:physics:phyx-mini-0364:phyxminiproblems-problemphyxmini0364-pressureinpascals, def:physics:phyx-mini-0364:phyxminiproblems-problemphyxmini0364-densityinkilogramspercubicmeter, def:physics:phyx-mini-0364:phyxminiproblems-problemphyxmini0364-accelerationinmeterspersecondsquared, def:physics:phyx-mini-0364:phyxminiproblems-problemphyxmini0364-closedtoppipesetup}
  Standard textbook calibrations implicit in the multiple-choice numerical answer: atmospheric pressure is approximated by 1.00 * 10\textasciicircum{}5 Pa, water by 1000 kg/m\textasciicircum{}3, and gravity by 9.8 m/s\textasciicircum{}2. None mentions the requested trapped-air length
\end{definition}
\begin{proof}
  This is the declared model definition of Uses Standard Atmosphere Water And Gravity.
\end{proof}
\begin{definition}[Has Physical Pipe Parameters]
  \label{def:physics:phyx-mini-0364:phyxminiproblems-problemphyxmini0364-hasphysicalpipeparameters}
  \lean{PhyXMiniProblems.ProblemPhyXMini0364.HasPhysicalPipeParameters}
  \uses{def:physics:phyx-mini-0364:phyxminiproblems-problemphyxmini0364-lengthinmeters, def:physics:phyx-mini-0364:phyxminiproblems-problemphyxmini0364-areainsquaremeters, def:physics:phyx-mini-0364:phyxminiproblems-problemphyxmini0364-volumeincubicmeters, def:physics:phyx-mini-0364:phyxminiproblems-problemphyxmini0364-pressureinpascals, def:physics:phyx-mini-0364:phyxminiproblems-problemphyxmini0364-densityinkilogramspercubicmeter, def:physics:phyx-mini-0364:phyxminiproblems-problemphyxmini0364-accelerationinmeterspersecondsquared, def:physics:phyx-mini-0364:phyxminiproblems-problemphyxmini0364-closedtoppipesetup}
  Positivity and geometric ordering for the two physical states. The final air column is shorter than the initially full pipe, selecting compression without assigning the final length any answer-choice value
\end{definition}
\begin{proof}
  This is the declared model definition of Has Physical Pipe Parameters.
\end{proof}
\begin{definition}[Satisfies Closed Top Pipe Laws]
  \label{def:physics:phyx-mini-0364:phyxminiproblems-problemphyxmini0364-satisfiesclosedtoppipelaws}
  \lean{PhyXMiniProblems.ProblemPhyXMini0364.SatisfiesClosedTopPipeLaws}
  \uses{def:physics:phyx-mini-0364:phyxminiproblems-problemphyxmini0364-lengthreadout, def:physics:phyx-mini-0364:phyxminiproblems-problemphyxmini0364-areareadout, def:physics:phyx-mini-0364:phyxminiproblems-problemphyxmini0364-volumereadout, def:physics:phyx-mini-0364:phyxminiproblems-problemphyxmini0364-pressurereadout, def:physics:phyx-mini-0364:phyxminiproblems-problemphyxmini0364-densityreadout, def:physics:phyx-mini-0364:phyxminiproblems-problemphyxmini0364-accelerationreadout, def:physics:phyx-mini-0364:phyxminiproblems-problemphyxmini0364-pipestate, def:physics:phyx-mini-0364:phyxminiproblems-problemphyxmini0364-closedtoppipesetup}
  The governing relations for the experiment: constant cross section gives V = A * length in each state; before immersion the freely connected air is at atmospheric pressure; the slowly compressed fixed gas obeys Boyle's law P\_i V\_i = P\_f V\_f; at the final internal interface, hydrostatics gives P\_f = P\_atm + rho * g * depth. The equations are stated in every coherent selection of base units. They do not give the final length a numerical value and do not mention an answer choice
\end{definition}
\begin{proof}
  This is the declared model definition of Satisfies Closed Top Pipe Laws.
\end{proof}
\begin{definition}[Answer Choice]
  \label{def:physics:phyx-mini-0364:phyxminiproblems-problemphyxmini0364-answerchoice}
  \lean{PhyXMiniProblems.ProblemPhyXMini0364.AnswerChoice}
  Labels of the four displayed answer choices
\end{definition}
\begin{proof}
  This is the declared model definition of Answer Choice.
\end{proof}
\begin{definition}[displayed Length In Meters]
  \label{def:physics:phyx-mini-0364:phyxminiproblems-problemphyxmini0364-displayedlengthinmeters}
  \lean{PhyXMiniProblems.ProblemPhyXMini0364.displayedLengthInMeters}
  \uses{def:physics:phyx-mini-0364:phyxminiproblems-problemphyxmini0364-answerchoice}
  Trapped-air length in meters printed beside each answer label
\end{definition}
\begin{proof}
  This is the declared model definition of displayed Length In Meters.
\end{proof}
\begin{definition}[recorded Dataset Answer]
  \label{def:physics:phyx-mini-0364:phyxminiproblems-problemphyxmini0364-recordeddatasetanswer}
  \lean{PhyXMiniProblems.ProblemPhyXMini0364.recordedDatasetAnswer}
  \uses{def:physics:phyx-mini-0364:phyxminiproblems-problemphyxmini0364-answerchoice}
  Dataset answer metadata, deliberately not used as a premise
\end{definition}
\begin{proof}
  This is the declared model definition of recorded Dataset Answer.
\end{proof}
\begin{definition}[Rounds To Nearest Tenth Meter]
  \label{def:physics:phyx-mini-0364:phyxminiproblems-problemphyxmini0364-roundstonearesttenthmeter}
  \lean{PhyXMiniProblems.ProblemPhyXMini0364.RoundsToNearestTenthMeter}
  \uses{def:physics:phyx-mini-0364:phyxminiproblems-problemphyxmini0364-lengthquantity, def:physics:phyx-mini-0364:phyxminiproblems-problemphyxmini0364-lengthinmeters}
  A physical length rounds to a displayed value to the nearest tenth meter
\end{definition}
\begin{proof}
  This is the declared model definition of Rounds To Nearest Tenth Meter.
\end{proof}
\begin{definition}[Matches Displayed Trapped Air Length]
  \label{def:physics:phyx-mini-0364:phyxminiproblems-problemphyxmini0364-matchesdisplayedtrappedairlength}
  \lean{PhyXMiniProblems.ProblemPhyXMini0364.MatchesDisplayedTrappedAirLength}
  \uses{def:physics:phyx-mini-0364:phyxminiproblems-problemphyxmini0364-closedtoppipesetup, def:physics:phyx-mini-0364:phyxminiproblems-problemphyxmini0364-answerchoice, def:physics:phyx-mini-0364:phyxminiproblems-problemphyxmini0364-displayedlengthinmeters, def:physics:phyx-mini-0364:phyxminiproblems-problemphyxmini0364-roundstonearesttenthmeter}
  A displayed choice agrees with the final trapped-air length
\end{definition}
\begin{proof}
  This is the declared model definition of Matches Displayed Trapped Air Length.
\end{proof}
\begin{definition}[Is Unique Matching Displayed Length]
  \label{def:physics:phyx-mini-0364:phyxminiproblems-problemphyxmini0364-isuniquematchingdisplayedlength}
  \lean{PhyXMiniProblems.ProblemPhyXMini0364.IsUniqueMatchingDisplayedLength}
  \uses{def:physics:phyx-mini-0364:phyxminiproblems-problemphyxmini0364-closedtoppipesetup, def:physics:phyx-mini-0364:phyxminiproblems-problemphyxmini0364-answerchoice, def:physics:phyx-mini-0364:phyxminiproblems-problemphyxmini0364-matchesdisplayedtrappedairlength}
  Exactly one displayed choice agrees with the modeled final air length
\end{definition}
\begin{proof}
  This is the declared model definition of Is Unique Matching Displayed Length.
\end{proof}
\begin{lemma}[final Air Length satisfies quadratic]
  \label{lem:physics:phyx-mini-0364:phyxminiproblems-problemphyxmini0364-finalairlength-satisfies-quadratic}
  \lean{PhyXMiniProblems.ProblemPhyXMini0364.finalAirLength_satisfies_quadratic}
  \uses{def:physics:phyx-mini-0364:phyxminiproblems-problemphyxmini0364-lengthinmeters, def:physics:phyx-mini-0364:phyxminiproblems-problemphyxmini0364-closedtoppipesetup, def:physics:phyx-mini-0364:phyxminiproblems-problemphyxmini0364-matchesprimaryfigure, def:physics:phyx-mini-0364:phyxminiproblems-problemphyxmini0364-matchesclosedtoppipescenario, def:physics:phyx-mini-0364:phyxminiproblems-problemphyxmini0364-usesstandardatmospherewaterandgravity, def:physics:phyx-mini-0364:phyxminiproblems-problemphyxmini0364-hasphysicalpipeparameters, def:physics:phyx-mini-0364:phyxminiproblems-problemphyxmini0364-satisfiesclosedtoppipelaws}
  Eliminating the two air volumes and the final pressure from pipe geometry, Boyle's law, hydrostatics, and the final surface geometry gives 49 * L\textasciicircum{}2 + 500 * L - 1500 = 0 for the final trapped-air length L in meters. This is a derived relation, not a premise or definition of the final state
\end{lemma}
\begin{proof}
  The conclusion follows from the governing relations and figure readouts named in the statement of final Air Length satisfies quadratic.
\end{proof}
% --- Archon declaration topology end ---
% --- Archon physics formalization source end ---
